\section{Neurosymbolic Integration}\label{sec:nesy}

This section is where the three pillars meet. The device-resident substrate (\Cref{sec:datalog}) keeps the symbolic state on the GPU; verified knowledge compilation (\Cref{sec:prob}) turns a network's outputs into a sound differentiable circuit; and zero-copy coupling carries gradients back into the neural optimizer without leaving the device. Together they let neural perception and logical reasoning train jointly, end to end, with no host--device transfers in the loop.

\subsection{Neural predicates}\label{subsec:neural-pred}

In xlog a neural network \emph{is} a predicate. The \texttt{nn/4} declaration embeds a network directly into the program:

\begin{lstlisting}[style=xlogstyle,
                   caption={Neural predicate declaration (\texttt{nn/4}).},
                   label={lst:nn-decl}]
nn(network_name, [input_vars], output_var,
   [output_labels]) :: predicate(args).
\end{lstlisting}

This is not a foreign-function escape hatch: it states that evaluating \texttt{predicate(args)} requires a forward pass through the named network, whose softmax over the label set becomes a distribution over probabilistic facts feeding the knowledge-compilation pipeline of \Cref{sec:prob}. On the Python side a PyTorch module is registered against the predicate:

\begin{lstlisting}[style=pythonstyle,
                   caption={Registering a PyTorch module against a neural predicate.},
                   label={lst:nn-register}]
program.register_network("mnist_net", net, optimizer)
\end{lstlisting}

Registration is typed, not by name. A caller may state the arity of the network's predicate, and that claim is checked against every \texttt{nn/4} declaration bound to the network in the program and rejected on mismatch: the program's own declaration is the authority. Per-argument catalog sort identifiers may accompany the arity (one per declared argument, refused without an arity or at the wrong length) and an opaque content hash may identify the registered network, so a retrained one mints a new identity when it is registered again. The engine interprets neither; both are carried for consumers, which read them back through a queryable metadata view together with each declaration's predicate, predicate arity, input arity, and label set. A downstream rule learner can therefore check a candidate against what the program declares instead of against a naming convention---the property \Cref{subsec:engine-mode} relies on.

The dataflow runs in two directions. \textbf{Forward:} when evaluation reaches an \texttt{nn/4}-backed predicate, the registered module runs, and its softmax vector is decomposed into weighted probabilistic facts, one per label, and fed into the knowledge-compilation pipeline of \Cref{sec:prob}, whose circuit cache carries across iterations. \textbf{Backward:} after the circuit computes the forward log-probability and backward gradients via level-parallel kernels, the per-leaf gradient buffers are exported as DLPack tensors and injected into PyTorch's autograd graph, and the optimizer updates the network as usual. The GIL is released during GPU work, so CUDA execution does not block the interpreter. A neural predicate is not confined to a query head: it may also appear inside rule bodies, so that whole rules, not just a single perceptual predicate, become the trainable object (\Cref{subsec:trainable-bodies}).

\subsection{End-to-end training}

The canonical example is MNIST addition (the DeepProbLog benchmark): given two digit images, predict their sum, with supervision only on sums: the network never sees individual digit labels. The entire reasoning structure is two rules:

\begin{lstlisting}[style=xlogstyle,
                   caption={MNIST addition: a CNN-backed digit predicate and an addition rule.},
                   label={lst:mnist-prog}]
nn(mnist_net, [X], Y, [0,1,2,3,4,5,6,7,8,9])
    :: digit(X, Y).
addition(X, Y, Z) :-
    digit(X, D1), digit(Y, D2), Z is D1 + D2.
\end{lstlisting}

The \texttt{nn/4} declaration connects the CNN to \texttt{digit/2}; the addition rule computes every consistent decomposition of a sum through probabilistic marginalization. Each query \texttt{addition(i, j, s)} trains by minimizing $-\log P(\texttt{addition}(i, j, s))$ under the compiled circuit. The driver compiles the program, registers the network, and runs the loop, which batches queries sharing a circuit template and runs forward/backward over the cached circuit:

\begin{lstlisting}[style=pythonstyle,
                   caption={Python driver: compile, register, train.},
                   label={lst:mnist-train}]
program = pyxlog.Program.compile(SRC)
program.register_network("mnist_net", net, optimizer)
program.add_tensor_source("train", train_images)
history = pyxlog.train_model_tensor(
    program, queries,
    epochs=epochs, batch_size=batch_size)
\end{lstlisting}

\begin{figure*}[tbp]
  \centering
  \begin{tikzpicture}[x=1mm,y=1mm,
      ring/.style={anchor=north west,text width=26mm,
                   inner xsep=2mm,inner ysep=1.8mm,minimum height=14.5mm}]
    % ---- forward row ----
    \node[fneural,ring] (net) at (0,0)
      {neural network\\[0.4mm]\fsubt{PyTorch forward}};
    \node[fproc,ring] (ground) at (54,0)
      {symbolic grounding\\[0.4mm]\fsubt{rules $\cdot$ provenance}};
    \node[fproc,ring] (circ) at (96,0)
      {verified circuit\\[0.4mm]\fsubt{proven once $\cdot$ cached}};
    \node[fneural,ring] (loss) at (138,0)
      {loss $-\log P$\\[0.4mm]\fsubt{per-query NLL}};
    \node[circle,draw=figaccent,fill=white,line width=0.7pt,inner sep=0.4pt,
          font=\tiny] at (circ.north west) {\checkmark};
    % ---- backward row ----
    \node[fproc,ring] (back) at (138,-27)
      {circuit backward\\[0.4mm]\fsubt{level-parallel adjoint}};
    \node[fneural,ring] (opt) at (0,-27)
      {optimizer step\\[0.4mm]\fsubt{PyTorch autograd}};
    % ---- ring edges ----
    \draw[farrow] (net) -- node[flabel,above=0.4mm,text width=21mm]
      {softmax $\to$ weighted facts} (ground);
    \draw[farrow] (ground) -- (circ);
    \draw[farrow] (circ) -- node[flabel,above=0.4mm,text width=11mm]
      {forward $\log P$} (loss);
    \draw[fback] (loss) -- node[flabel,anchor=east,xshift=-1mm] {backward} (back);
    \draw[fback] (back) -- node[flabel,above=0.4mm]
      {per-leaf gradients exported via DLPack (zero-copy)} (opt);
    \draw[fback] (opt) -- node[flabel,anchor=west,xshift=1mm]
      {update $\theta$} (net);
    \node[flabel] at (84,-20)
      {one training iteration --- circuit \& proof reused across iterations};
    % ---- GPU device plane around the whole loop ----
    \node[flabel] (foot) at (84,-45.5)
      {zero host--device transfers in the loop (byte-accounted)};
    \begin{scope}[on background layer]
      \node[fplane,fit=(net)(loss)(back)(opt)(foot),inner sep=3mm] (plane) {};
    \end{scope}
    \node[fplanetitle,anchor=west] at ([xshift=4mm]plane.north west)
      {\planetitle{gpu device-resident}};
    % ---- legend ----
    \draw[figgreen,line width=0.9pt,fill=white,rounded corners=1pt]
      (38,-56.5) rectangle +(3,3);
    \node[fsub,anchor=west] at (42.5,-55) {PyTorch};
    \draw[figblue,line width=0.9pt,fill=white,rounded corners=1pt]
      (62,-56.5) rectangle +(3,3);
    \node[fsub,anchor=west] at (66.5,-55) {xlog substrate};
    \draw[figaccent,line width=0.9pt,fill=white,rounded corners=1pt]
      (96,-56.5) rectangle +(3,3);
    \node[fsub,anchor=west] at (100.5,-55) {verified (\Cref{fig:kc})};
  \end{tikzpicture}
  \caption{One end-to-end neurosymbolic training iteration. The entire
  loop---neural forward pass, symbolic grounding, circuit evaluation, loss,
  and gradient return---executes on one CUDA device (shaded plane) with zero
  host--device data-plane transfers. The circuit is compiled and proven
  correct once (\checkmark, \Cref{fig:kc}), then reused across iterations,
  and gradients re-enter PyTorch's autograd graph zero-copy through DLPack.}
  \label{fig:training}
\end{figure*}

The decisive performance property (\Cref{fig:training}) is that compilation and verification happen once, on the first forward pass. Every later iteration reuses the cached circuit, executing only the weight update and the level-parallel forward/backward kernels. This is the mechanism behind the $2.74\times$ training speedup; the loss reduction across seeds confirms that gradients flow from circuit evaluation through the logic program back to the CNN, and the full timing breakdown is reported in \Cref{sec:eval}.

\subsection{Zero-copy interoperability}\label{subsec:interop}

A GPU-native engine is only useful if results reach the frameworks researchers already use, without copies that would reintroduce the transfer wall. xlog exposes its device-resident state through four standard interfaces.

\textbf{DLPack.} Each relation column or gradient buffer becomes a managed DLPack~\cite{dlpack} tensor whose GPU memory is reference-counted and freed only when every consumer releases it; on import, xlog validates device, dtype, contiguity, and alignment, optionally against an expected schema. In Python the protocol surfaces as a capsule, so \texttt{torch.from\_dlpack} yields a live CUDA tensor backed by the very memory xlog wrote.

\textbf{Arrow.} For columnar tools (Polars, DuckDB, cuDF), relations export through the Arrow~\cite{arrow} C Device Data Interface; symbol columns export as Arrow dictionary arrays, so a consumer reads symbolic columns as native string dictionaries while xlog keeps the compact integer representation.

\textbf{Python bindings.} The \texttt{pyxlog} extension (PyO3~\cite{pyo3}) exposes compilation, evaluation, training, and result extraction; tensor-valued results are returned as DLPack capsules, and long-running GPU operations release the GIL so data-loader threads proceed concurrently.

\textbf{PyTorch autograd.} The circuit behaves as a differentiable function inside PyTorch: the network's grad-enabled forward feeds the circuit, the scalar loss is exported via DLPack and re-imported, and a batched path stacks inputs into one forward pass and runs a single backward, so gradients flow from the logic layer back to \texttt{nn.Module} parameters with no custom autograd \texttt{Function}.

\subsection{Trainable rule bodies}\label{subsec:trainable-bodies}

Placing a neural predicate in a rule body lifts xlog from learning a single perceptual predicate to learning whole rules, while the pipeline stays device-resident with no host round-trips.

\textbf{Mixed bodies.} A trainable body may interleave neural-evaluated atoms with ordinary symbolic literals. The symbolic literals act as hard join conditions: they gate which groundings can fire but carry no probability mass or gradient. The head probability therefore factors as a $0/1$ indicator of the symbolic condition times the neural probability scaled by a learned per-rule weight~$\sigma(w)$, and gradients flow only through the network outputs and the rule weight, never through the ground facts.

\textbf{Existential joins.} When an existential (non-head) body variable is a neural predicate's input, xlog grounds that predicate over the \emph{real} join domain \emph{inside} the circuit instead of stripping the join as a pre-filter: the neural occurrence expands into one circuit leaf per domain event, and provenance OR-aggregates $\exists\,e.\ \mathrm{neural}(e)\wedge\mathrm{join}(e,\mathit{head})$ at each head binding. This makes ``does supporting evidence exist?'' rules learnable, reasoning over a whole domain of events instead of only per-instance classification.

\textbf{Joint multi-rule mixtures.} Several same-head trainable clauses combine as a noisy-OR mixture, $P(\mathit{head}_i)=1-\prod_k\bigl(1-r_k[i]\,m_k[i]\,\sigma(w_k)\bigr)$. Here $r_k[i]$ is a candidate's relational eligibility, $\sigma(w_k)$ its learned rule weight, and $m_k[i]$ an optional graded confidence in $[0,1]$ carried per binding, for example a fact's confidence trajectory over time. Rule learning thus moves from binary rule eligibility to a continuous, differentiable weighting of competing and cooperating rules.

\textbf{Learned body gates.} Inside that mixture a candidate may also carry a neural conjunct over a per-binding entity feature $\phi(x)$: its eligibility becomes its relational grounding \emph{and} a straight-through-thresholded gate $g_\theta(\phi(x))\ge\tau$. A gated clause then competes against ungated ones on the same noisy-OR terms, and $g_\theta$ trains jointly with the rule weights under the same held-out selector. By default $\phi$ is detached where it enters the engine: gradient reaches $\theta$ and the rule weight but stops at the neural head, leaving the feature producer frozen. Training the producer too is an explicit opt-in per candidate: one flag leaves $\phi$ attached, so gradient flows on into the backbone that computed it. Only the autograd linkage changes; the values uploaded are identical either way.

These trainable-rule paths share the verified circuit and the zero-copy substrate of the simpler predicate case, and extend to recursive programs now that provenance converges on two-sided recursive strongly-connected components. \Cref{sec:ecinduction} carries the surface onto a public benchmark, not a demonstration: the perceptual predicate inside an induced Event-Calculus rule body is trained through symbolic credit alone, cross-validated under pre-registered gates. \Cref{sec:limits} states what that evidence does and does not support.

\subsection{Engine-mode credit and candidate selection}\label{subsec:engine-mode}

Which of the candidate bodies the engine enumerates should the learner keep, and on what evidence? Learning a rule body is not only a gradient problem: it is also a decision about which candidate to keep, and the two need different evidence. Engine-mode training separates them. The engine already enumerates the candidate space (the well-formed pairs of body relations over the program's own binary relations) and already holds the join extensions those candidates read, so training scores that space directly instead of a reconstruction of it.

A candidate's score on a fact is a noisy-OR over the witnesses the engine reports for that fact: a neural body relation contributes a probability per witness, read at the fact's own label, while a purely relational one contributes a fixed $\{0,1\}$ cover. Candidates combine as a softmax over their logits, and the resulting negative log-likelihood is a single autograd graph, so one backward pass updates the candidate logits and the network behind the probabilities together.

\textbf{A worked instance.} Take one head fact and two candidates for its body, one relational and one neural. The relational candidate covers the fact or does not, so its score there is $1$ or $0$. The neural candidate's score is the noisy-OR over the witnesses the engine reports for that fact, each read at the fact's own label: two witnesses at probability $0.5$ give $1-0.5\cdot0.5=0.75$, above a relational candidate that does not cover the fact and below one that does. Those are training scores, and the arbiter never reads them. It rescores both on the held-out facts by the same witness semantics, drops either whose held-out score falls below the fit gate, and keeps the relational one when the two land within the score quantum; if that narrowing leaves two relational candidates the data cannot separate, it abstains rather than pick one.

The candidate pool is validated against the typed registry of \Cref{subsec:neural-pred} instead of being trusted by name: a neural relation whose declared arity the credit has no witness semantics for is refused at registration, and a pool that filtering empties reports its per-filter counts instead of silently shrinking. And because raw engine constants index the caller's feature rows directly, that identity is accepted only when the witness domain is exactly the dense row range; a misalignment that would stay in bounds while gathering another event's probability is refused, not guessed at.

\textbf{Selection on held-out generalization.} Training credit cannot distinguish a crisp but coincidental rule from a soft but correct one, even in principle; generalization can. Selection therefore never reads the trained weight. The facts are split into folds; each fold trains on the rest and rescores every candidate on the held-out facts by its own witness semantics; the arbiter ranks the fold-averaged scores. A \emph{fit gate} then drops any candidate whose held-out score falls below a threshold, since a rule that cannot fit held-out data is not a rule.

Ties within the score quantum break toward the relational candidate on Occam grounds, but only when that narrowing leaves exactly one: preferring a relational explanation over a neural one is licensed, choosing among relational duplicates the data cannot separate is not. When neither step is decisive the arbiter returns no rule and says why; abstention is an outcome of the selector, not a failure of it. The held-out score may be accuracy or F1, and the choice matters: at a rare positive class accuracy sits on the all-negative base-rate plateau, where a candidate that predicts nothing outranks the best real detector, which is why the benchmark protocol of \Cref{sec:ecinduction} selects on held-out F1 instead.

\textbf{Masked witnesses and withheld evidence.} A per-witness mask lets evidence be withheld rather than denied: a masked witness leaves the candidate's index entirely, contributing zero credit and zero gradient instead of a coerced false, and the facts it thereby leaves uncertain are excluded from the held-out score and counted against a coverage gate that runs before the fit gate, so no candidate is killed by facts it was never allowed to see. The channel rides the accuracy axis only and is refused up front beside the F1 holdout, whose measurability derives from raw labels a mask can hide.

\textbf{Acceptance without retraining.} The same machinery also runs as an acceptance instrument: given an externally trained detector, a frozen entry point scores every enumerated candidate against it with no optimizer, no gradient and no folds and applies exactly the arbiter's gates, so where the cross-validated path asks whether a detector can be trained for a rule, this one asks whether a rule is right given the detector one already has. A training-mode module is refused rather than quietly switched, since batch-norm statistics and dropout would de-determinize the bit-identical scores the entry point exists to provide.

\subsection{Term embeddings and differentiable ILP}

\textbf{Term embeddings} attach dense, optionally trainable vectors to logical symbols. Where \texttt{nn/4} maps perception to discrete facts, embeddings give logical entities a continuous representation that participates in neural computation, with gradients flowing through the embedding lookup: knowledge-graph entity representations, for instance, trained jointly with neural perception.

\textbf{Differentiable ILP} learns first-order rules instead of network weights. Candidate rules are represented as sparse bitmasks over the ground-atom space, credit assignment runs entirely on-device via scatter--gather, and a promotion pipeline decides which rules enter the program.

Four kinds of judgement sit on these two paths, and \Cref{tab:gates} separates them. Convergence and the supply of explicit held-out examples are preconditions checked before any gate runs, not gates: with neither held-out positives nor held-out negatives supplied, the pipeline returns for manual review before a gate is evaluated, so the always-run gates cannot promote on their own.

The benchmark runs of \Cref{sec:ecinduction} exercise the last row only. Its permutation-null threshold is not an extra gate but the holdout arbiter's own fit threshold, derived per fold from label permutations instead of fixed at a constant; the promotion gates above it belong to the differentiable-ILP path and are not on that route.

\begin{table}[htbp]
  \centering
  \caption{What judges what, on which path: the promotion pipeline of differentiable ILP or the holdout arbiter of \Cref{subsec:engine-mode}. Convergence and supplied held-out examples are preconditions, not gates: without them the pipeline returns for manual review before any row below runs.}
  \label{tab:gates}
  \footnotesize
  \setlength{\tabcolsep}{3pt}
  \begin{tabularx}{\columnwidth}{@{}>{\raggedright\arraybackslash}p{0.25\columnwidth}Ll@{}}
    \toprule
    \textbf{Judgement} & \textbf{What it judges} & \textbf{Path} \\
    \midrule
    Six always-run gates & the rule derives every training positive and no training negative, passes the novel-fact audit, loses no protected fact, meets a held-out F1 threshold, and names only relations with typed schemas & dILP \\
    \addlinespace
    Two held-out gates & the same rule on the supplied held-out positives and negatives & dILP \\
    \addlinespace
    Ambiguity scan & which other candidates also explain the data; reported, never decisive & dILP \\
    \addlinespace
    Permutation-null fit threshold & whether a candidate generalizes better than shuffled labels do on its own fold & arbiter \\
    \bottomrule
  \end{tabularx}
\end{table}

Both paths ride the device-resident bridge of \Cref{subsec:interop}. Where the trainable rule bodies of \Cref{subsec:trainable-bodies} learn rule \emph{weights} and per-binding confidences, differentiable ILP learns rule \emph{structure}; the two are complementary. This sparse, GPU-resident formulation avoids the cubic blowup of dense rule materialization. Unlike the trainable-body surface, this path is not exercised on the benchmark of \Cref{sec:ecinduction}, and it is a beta subsystem (\Cref{sec:limits}).
